% !TEX program = pdflatex
% vibe-summary.tex
%
% Self-contained English summary of the Master's thesis
%   "MẶT BLW VÀ BÀI TOÁN BJ\"ORLING" (BLW surfaces and the Björling problem)
%   Huỳnh Kỳ Anh, Master's thesis in Geometry and Topology,
%   Hue University of Education, 2006.
%   Advisor: Assoc. Prof. Đoàn Thế Hiếu; with the assistance of
%   Prof. José A. Gálvez.
%
% This file is an independent, standard-LaTeX rendition (no dependency
% on the thesis's custom `main.sty`/`ktv-maths` machinery) prepared so
% that the mathematical content of the (Vietnamese) thesis can be read
% and typeset in English. It is a SUMMARY, not a translation: proofs
% are sketched, not reproduced in full detail; all definitions and
% theorem statements are given precisely.

\documentclass[11pt,a4paper]{article}

\usepackage[utf8]{inputenc}
\usepackage[T1]{fontenc}

% A handful of Vietnamese proper nouns in the front matter (the author's
% and advisor's names, and the original Vietnamese subtitle) use
% precomposed Unicode letters with stacked diacritics that inputenc's
% default utf8 table does not know. Teach it about them explicitly,
% using the LaTeX kernel's own accent commands (T1 fontenc, loaded
% above, already has a native glyph for \DJ, so \DH{} is exact; the
% double-accented letters are built by stacking two kernel accents,
% which pdfTeX renders without needing any extra Vietnamese-specific
% package such as vntex).
\DeclareUnicodeCharacter{0110}{\DJ}       % Đ (D with stroke)
\DeclareUnicodeCharacter{1EB6}{\d{\u A}}  % Ặ (A + breve + dot below)
\DeclareUnicodeCharacter{1EB7}{\d{\u a}}  % ặ (a + breve + dot below)
\DeclareUnicodeCharacter{1EBF}{\'{\^e}}   % ế (e + circumflex + acute)
\DeclareUnicodeCharacter{1EF3}{\`y}       % ỳ (y + grave)

\usepackage{amsmath,amssymb,amsthm,mathtools}
\usepackage{geometry}
\usepackage{enumitem}
\usepackage{hyperref}
\usepackage{bm}

\geometry{margin=1in}

% =====================================================================
% Notation macros (standard LaTeX, self-contained)
% =====================================================================
\newcommand{\RR}{\mathbb{R}}
\newcommand{\CC}{\mathbb{C}}
\newcommand{\Cbar}{\overline{\mathbb{C}}}           % Riemann sphere
\newcommand{\Hyp}{\mathbb{H}^3}                      % hyperbolic 3-space
\newcommand{\Lor}{\mathbb{L}^4}                      % Lorentz 4-space
\newcommand{\deS}{\mathbb{S}^3_1}                    % de Sitter space
\newcommand{\LCone}{\mathbb{N}^3_+}                  % future light cone
\newcommand{\Herm}{\operatorname{Herm}(2)}
\newcommand{\SLtwo}{\operatorname{SL}(2,\CC)}
\newcommand{\ip}[2]{\langle #1,#2\rangle}            % Lorentz inner product
\newcommand{\cross}[2]{#1 \wedge #2}                 % Lorentz cross product
\newcommand{\dz}{\partial_z}
\newcommand{\dzb}{\partial_{\bar z}}
\newcommand{\sgn}{\operatorname{sgn}}

\theoremstyle{plain}
\newtheorem{theorem}{Theorem}
\newtheorem{corollary}[theorem]{Corollary}
\newtheorem{lemma}[theorem]{Lemma}
\newtheorem{proposition}[theorem]{Proposition}
\theoremstyle{definition}
\newtheorem{definition}[theorem]{Definition}
\newtheorem{problem}[theorem]{Problem}
\theoremstyle{remark}
\newtheorem{remark}[theorem]{Remark}

\title{\textbf{BLW Surfaces and the Björling Problem}\\[2mm]
\large English summary of the Master's thesis\\
\textit{Mặt BLW và Bài toán Björling}}
\author{Huỳnh Kỳ Anh\\[1mm]
\small Master's Thesis in Geometry and Topology (code 60 46 10)\\
\small Hue University of Education, 2006\\
\small Advisor: Assoc.\ Prof.\ Đoàn Thế Hiếu\\
\small with the assistance of Prof.\ José A.\ Gálvez\\[3mm]
\small \emph{(Summary prepared in English from the original Vietnamese thesis)}}
\date{}

\begin{document}
\maketitle

\begin{abstract}
The classical \emph{Björling problem}, posed by E.\,G.\ Björling in 1844 and
solved by H.\,A.\ Schwarz in 1890, asks for the minimal surface in $\RR^3$
containing a prescribed analytic curve with a prescribed unit normal field
along it. This thesis poses and solves the analogous problem for the class
of \emph{BLW surfaces} (``Linear Weingarten surfaces of Bryant type'') in
hyperbolic $3$-space $\Hyp$: surfaces whose mean curvature $H$ and (extrinsic)
Gauss curvature $K_I$ satisfy the linear relation $2a(H-1)+bK_I=0$.
This class contains the $\mathtt{CMC1}$ surfaces (constant mean curvature
one) of Gálvez--Mira \cite{GMM1} ($b=0$) and the flat surfaces of
Gálvez--Mira \cite{GMM5} ($a=0$) as special cases. The main results are:
(1) a full existence-and-uniqueness theorem for the Björling problem for BLW
surfaces, together with an explicit closed-form representation of the
solution (a ``Björling representation'' that does not go through the
conformal/harmonic representation of \cite{GMM2}); and (2) a theorem
determining exactly when the hyperbolic Gauss map of a BLW surface is
meromorphic versus anti-meromorphic, generalizing R.\ Bryant's classical
fact that the hyperbolic Gauss map of a regular surface in $\Hyp$ is
meromorphic iff the surface has $H\equiv 1$ (i.e., is $\mathtt{CMC1}$).
\end{abstract}

\tableofcontents

% =====================================================================
\section{Introduction}
% =====================================================================

The \emph{classical Björling problem} was introduced by E.\,G.\ Björling in
1844 and solved by H.\,A.\ Schwarz in 1890: given an analytic curve
$\beta\subset\RR^3$ with an analytic unit tangent, find the minimal surface
containing $\beta$ as a geodesic. A key structural consequence of the
existence-and-uniqueness of the solution is that a minimal surface $\Sigma$
is symmetric with respect to any plane orthogonal to $\Sigma$.

A natural generalization is: given a class $\mathfrak B$ of surfaces
(e.g.\ $\mathtt{CMC1}$, $\mathtt{ELW}$, $\mathtt{flat}$, \dots), a curve
$\beta$, and a unit normal field $V$ along $\beta$, find a surface
$\Sigma\in\mathfrak B$ containing $\beta$ with prescribed unit normal $V$
along it. Appendix~A of the thesis lists the Björling-type problems that
had been solved before this work; in every known case the solution is
unique and yields a symmetry principle for the surface class in question.
Solving the Björling problem for a class also produces many explicit
examples and has led to classification results — e.g.\ the classification
of $\mathtt{CMC1}$ surfaces of cylindrical topology \cite{GMM1}, and of
flat fronts with embedded isolated singularities \cite{GMM5}.

\medskip
Building on the solutions for $\mathtt{CMC1}$ surfaces \cite{GMM1} and
regular flat surfaces \cite{GMM5}, this thesis poses and solves the
Björling problem for \textbf{BLW surfaces}: surfaces $\psi:S\to\Hyp$ whose
mean curvature $H$ and Gauss curvature $K_I=K-1$ (here $K$ is the
Gauss--Kronecker/extrinsic curvature) satisfy
\begin{equation}\label{eq:BLW}
  2a(H-1)+bK_I = 0 ,
\end{equation}
for real constants $a,b$ with $a+b\neq0$. The technique follows
\cite{GMM1}: to build the solution one needs a fact about the
meromorphicity of the hyperbolic Gauss map, originally due to
R.\ Bryant \cite{BRYANT}. This fact is partially generalized in
Theorem~\ref{thm:G} below, which is the key ingredient for solving the
Björling problem for BLW surfaces. (Open questions about further
generalizing Bryant's theorem are discussed in
Section~\ref{sec:conclusion}, Problem~\ref{open:G}.)

\subsection*{Outline}
Section~\ref{sec:prelim} collects the background material: the model of
hyperbolic space, the local theory of surfaces in $\Hyp$, Riemann
surfaces and (anti-)meromorphic maps, conformality, basic lemmas on BLW
surfaces, and the Cauchy problem for the Liouville equation.
Section~\ref{sec:BJ} states the Björling problem for BLW surfaces and the
two main theorems (Theorem~\ref{thm:BJ}, the existence-and-uniqueness
theorem, and Theorem~\ref{thm:G}, on meromorphicity of the hyperbolic
Gauss map), and sketches their proofs.
Section~\ref{sec:conclusion} summarizes the contributions and lists
several open problems suggested by the work.
Appendix~\ref{app:list} reproduces the thesis's table of previously solved
Björling-type problems.

% =====================================================================
\section{Preliminaries}\label{sec:prelim}
% =====================================================================

\subsection{The hyperbolic space $\Hyp$}\label{ssec:H3}

\emph{Lorentz $4$-space} $\Lor$ is $\RR^4$ equipped with the Lorentzian
pseudo-metric
\begin{equation}
  \ip{\cdot}{\cdot} = -dx_0^2+dx_1^2+dx_2^2+dx_3^2 .
\end{equation}
Two distinguished subsets are the \emph{de Sitter $3$-space}
$\deS=\{x\in\Lor : \ip xx=1\}$ and the \emph{future light cone}
$\LCone=\{x\in\Lor : \ip xx=0,\ x_0>0\}$. The set
\[
  \Hyp = \{x\in\Lor : \ip xx=-1,\ x_0>0\}
\]
is the \emph{Minkowski model of (3-dimensional) hyperbolic space}; its
tangent space at $p$ is $T_p\Hyp=\{x\in\Lor:\ip xp=0\}$.

Identifying $(x_0,x_1,x_2,x_3)\in\Lor$ with the Hermitian $2\times2$ matrix
$\begin{psmallmatrix}x_0+x_3 & x_1+ix_2\\ x_1-ix_2 & x_0-x_3\end{psmallmatrix}$
realizes $\Lor$ as $\Herm$, the space of Hermitian $2\times2$ matrices,
with $\ip mm=-\det(m)$; thus $\Hyp=\{m\in\Herm:\det m=1\}$. This
matrix model is used throughout the construction of solutions.

For $p,u,v\in\Lor$ the \emph{Lorentzian cross product} (``triple wedge
product'') is
\begin{equation}
  p\wedge u\wedge v = -\det\begin{pmatrix} -i&j&k&l\\ p_0&p_1&p_2&p_3\\
  u_0&u_1&u_2&u_3\\ v_0&v_1&v_2&v_3\end{pmatrix},
\end{equation}
$(i,j,k,l)$ the standard basis, characterized by
$\ip{p\wedge u\wedge v}{w}=\det(p,u,v,w)$ for all $w$. Three identities
proved directly from the determinant (identities \eqref{eq:wpbase}--\eqref{eq:wpb}
below; these are singled out by the author as small original computations
used repeatedly in the proofs) are, for $\ip uv=0$, $w=p\wedge u\wedge v$, $\alpha=\ip uu$,
$\beta=\ip vv$:
\begin{gather}
  \ip ww = -\alpha\beta\ip pp + \alpha\ip pv^2+\beta\ip pu^2, \label{eq:wpbase}\\
  \ip{u\wedge v}{u\wedge v} = \ip uu\ip vv - \ip uv^2
    \qquad (u,v,w\in T_p\Hyp), \label{eq:wpa}\\
  w\wedge(u\wedge v) = u\ip wv - v\ip wu \qquad (u,v,w\in T_p\Hyp). \label{eq:wpb}
\end{gather}
The Lorentz metric restricted to each $T_p\Hyp$ is positive definite, so
$\Hyp$ is Riemannian; by do Carmo \cite{CARMO} it is simply connected,
complete, of constant sectional curvature $-1$. Note also: if
$(\beta,v)\in\Hyp\times\deS$ with $\ip\beta v=0$, then $n=\beta+v\in\LCone$
and $n_0+n_3\ge0$.

\subsection{Local theory of surfaces in $\Hyp$}\label{ssec:local}

Let $S$ be a connected, oriented $2$-manifold (a \emph{surface}) and
$\psi:S\to\Hyp$. A point $z_0$ is \emph{regular} if $\mathrm{rank}(d\psi)_{z_0}=2$;
in a conformal parameter $z=s+it$ this is $\{\psi_s\wedge\psi_t\}_{z_0}\ne0$.
The \emph{unit normal field} $\eta:S\to\deS$ satisfies $\ip\psi\eta=0$,
$\ip{d\psi}\eta=0$, and (matching the orientation) at regular points
\begin{equation}
  \eta = \frac{\psi_s\wedge\psi_t}{\sqrt{\ip{\psi_s\wedge\psi_t}{\psi_s\wedge\psi_t}}}.
\end{equation}
At regular points there is an expansion
\begin{equation}\label{eq:eta-expansion}
  \eta_z = -\Delta\,\psi_z+\Upsilon\,\psi_{\bar z}, \qquad
  \eta_{\bar z} = \bar\Upsilon\,\psi_z-\Delta\,\psi_{\bar z},
\end{equation}
with $\Delta$ real-valued and $\Upsilon$ complex-valued; the mean
curvature and Gauss--Kronecker curvature are $H=\Delta$ and
$K=H^2-|\Upsilon|^2$. Both are invariant under orientation-preserving
reparametrization.

The \emph{hyperbolic Gauss map} of $\psi$ is $G=[\psi+\eta]:S\to\Cbar$,
\begin{equation}\label{eq:Gdef}
  G = [\psi+\eta] = \frac{N_1-iN_2}{N_0+N_3}, \qquad N=\psi+\eta,
\end{equation}
i.e.\ the south-pole stereographic projection of the ideal boundary
$\partial^2_\infty\Hyp\cong\Cbar$ onto the extended complex plane through the
center of $\partial^2_\infty\Hyp$ (with $G=\infty$ when $N_0+N_3=0$). See
Bryant \cite{BRYANT} and Gálvez--Mira \cite{GMM2,GMM1} for details.

\subsection{Riemann surfaces; meromorphic and anti-meromorphic maps}

A \emph{Riemann surface} is a $2$-manifold with an atlas of holomorphically
compatible charts. Every oriented Riemannian $2$-manifold $(X,\sigma)$
admits such an atlas in which $\sigma$ is conformal (isothermal
coordinates), so every oriented Riemannian surface may be regarded as a
Riemann surface. The Riemann sphere $\partial^2_\infty\Hyp\cong\Cbar$ is
an example.

A \emph{meromorphic} function $f:S\to\Cbar$ is a holomorphic map in the
manifold sense; equivalently, in an orientation-preserving conformal
parameter $z$,
\begin{equation}\label{eq:meromorphic}
  f_{\bar z} = 0.
\end{equation}
Analogously $f$ is \emph{anti-meromorphic} if $f_z=0$ in such parameters.
Reversing the orientation of $\Cbar$ interchanges meromorphic and
anti-meromorphic maps, so the two classes share properties such as
uniqueness theorems and the identity theorem.

\subsection{Conformality}

A regular map $f:(M,\sigma_1)\to(N,\ip\cdot\cdot)$ between Riemannian
manifolds is \emph{conformal} if $\sigma_1$ is a positive multiple of the
pulled-back metric $\sigma_2(u,v):=\ip{df(u)}{df(v)}$. When $\dim M=2$
and $z=s+it$ is $\sigma_1$-conformal, this is equivalent to
\begin{equation}\label{eq:fzfz}
  4\ip{f_z}{f_z} = \ip{f_s}{f_s}-\ip{f_t}{f_t}-2i\ip{f_s}{f_t}=0 .
\end{equation}
The notion extends to maps that are not everywhere regular and to targets
with only a pseudo-metric, by requiring \eqref{eq:fzfz} in every conformal
chart of $M$. A holomorphic map $f:D\subset\CC\to\CC$ is conformal at its
regular points (angle- and orientation-preserving); conversely a conformal
map $f:S\to\Cbar$ is holomorphic or anti-holomorphic, and vice versa.

\subsection{Basic lemmas on BLW surfaces}\label{ssec:BLWbasic}

\begin{definition}[BLW surface]\label{def:BLW}
A regular map $\psi:S\to\Hyp$ is a \textbf{BLW surface} (Linear Weingarten
surface of Bryant type) if there are real constants $a,b$ such that
\begin{equation}\label{eq:BLWdef}
  2a(H-1)+bK_I = 2a(H-1)+b(K-1) = 0,
\end{equation}
where $H$ is the mean curvature, $K$ the Gauss--Kronecker curvature, and
$K_I=K-1$ the (intrinsic) Gauss curvature of $\psi$.
\end{definition}

The degenerate case $a+b=0$ is treated in \cite{MAGIC} and is excluded
throughout ($a+b\neq0$ is assumed from here on). If $b=0$ then $a=\pm1$
and \eqref{eq:BLWdef} says $H\equiv1$: these are the $\mathtt{CMC1}$
surfaces, for which the Björling problem was solved in \cite{GMM1}. If
$a=0$, \eqref{eq:BLWdef} says $K\equiv1$, i.e.\ the surface is
\emph{flat}, as studied in \cite{GMM5}.

\begin{lemma}[{\cite[Lemma 1]{GMM2}}]\label{lem:sigma}
Let $\psi$ be a BLW surface satisfying \eqref{eq:BLWdef}. One may
normalize $|a+b|=1$, choosing the signs of $a,b$ so that
\begin{equation}
  \sigma = aI+bI\!I
\end{equation}
is a positive-definite metric on $S$, where $I=\ip{d\psi}{d\psi}$ and
$I\!I=\ip{d\psi}{-d\eta}$ are the first and second fundamental forms.
\end{lemma}
Throughout, $S$ is regarded as a Riemann surface for the conformal
structure induced by $\sigma$, and it is always assumed $|a+b|=1$ with
signs chosen as in Lemma~\ref{lem:sigma}.

\begin{remark}[assembled from {\cite[Thm.\,1, its proof, and Remark 3]{GMM2}}]\label{rmk:psipluseta}
Let $\psi:S\to\Hyp$ be BLW as in \eqref{eq:BLWdef}, $\eta$ its unit normal,
$z$ a $\sigma$-conformal parameter, and $2\chi$ the conformal factor of
$\sigma=2\chi|dz|^2$. Then:
\begin{enumerate}[label=(\arabic*)]
\item $\psi+\eta$ is conformal for $\sigma$, and either $\psi(S)$ is
  part of a \emph{horosphere}, or $\ip{d(\psi+\eta)}{d(\psi+\eta)}$ is a
  pseudo-metric of constant curvature $c=\tfrac a{a+b}$; in particular the
  hyperbolic Gauss map $G=[\psi+\eta]$ is conformal;
\item $z_0\in S$ is a singular point of $\psi+\eta$ iff $H(z_0)=K(z_0)=1$;
  hence every such singular point is an \emph{umbilic} point of $\psi$;
\item the Hopf differential $Q_I=\ip{\psi_z}{\psi_z}\,dz^2$ is holomorphic,
  so a BLW surface is either totally umbilic or has isolated umbilic
  points;
\item at regular points of $\psi+\eta$,
  \begin{equation}
    \psi = \frac{2a+b}{2(a+b)}(\psi+\eta)
      + \frac{a+b}{\chi(2H-K-1)}(\psi+\eta)_{z\bar z};
  \end{equation}
  if in addition $b\neq0$,
  \begin{equation}\label{eq:2H-K-1}
    \frac{2H-K-1}{a+b}
      = \frac1\chi\ip{(\psi+\eta)_z}{(\psi+\eta)_{\bar z}},
  \end{equation}
  and, with $\phi$ the conformal factor of $\ip{d(\psi+\eta)}{d(\psi+\eta)}$,
  \begin{equation}\label{eq:psi-recovered}
    \psi = \frac{2a+b}{2(a+b)}(\psi+\eta) + \frac2\phi(\psi+\eta)_{z\bar z}.
  \end{equation}
\end{enumerate}
\end{remark}
By \cite{GMM1}, formula \eqref{eq:psi-recovered} — the key device for
\emph{recovering} $\psi$ from $\psi+\eta$ — remains valid even when $b=0$.

\subsection{The Liouville equation}\label{ssec:Liouville}

For a planar domain $U\subset\CC$ and real constant $c$, the
\emph{Liouville equation} is
\begin{equation}\label{eq:Ldef}
  (\log\phi)_{z\bar z} = -\frac c2\phi.
\end{equation}
If $ds^2=\phi|dz|^2$ is a Riemannian metric on $U$, then $\phi$ solves
\eqref{eq:Ldef} iff $ds^2$ has constant curvature $c$. This link between
surface theory and \eqref{eq:Ldef} is what drives the whole construction:
solving the Björling problem for BLW surfaces reduces to solving a Cauchy
problem for \eqref{eq:Ldef} (with $c=a/(a+b)$, the curvature of
$\ip{d(\psi+\eta)}{\cdot}$ from Remark~\ref{rmk:psipluseta}). This
contrasts with the classical minimal-surface Björling problem in $\RR^3$,
which reduces merely to integrating functions with no real periods; here
one must solve a genuine PDE.

The relevant Cauchy problem is
\begin{equation}\label{eq:Lcauchy}
  \begin{cases}
    (\log\phi)_{z\bar z} = -\dfrac c2\phi,\\[1mm]
    \phi(s,0)=a(s), \qquad \phi_z(s,0)=b(s),
  \end{cases}
\end{equation}
with $a(s)$ real-analytic and non-negative and $b(s)$ complex-analytic
with $2\Re\{b(s)\}=a'(s)$. Write $\{f,z\}=\big(f''/f'\big)'-\tfrac12\big(f''/f'\big)^2$
for the Schwarzian derivative.

\begin{theorem}[{\cite[Thm.\,1]{GMM1}}, case $c\neq0$]\label{thm:GMM1-1}
Let $c\neq0$ and let $T(s):I\to\CC$ solve $\{T,s\}=\Upsilon(s)$, where
\[
  2\Upsilon(s)=2\Big(\frac{b(s)}{a(s)}\Big)'-\Big(\frac{b(s)}{a(s)}\Big)^2+c\,a(s).
\]
Set
\[
  \overline R(s)=\frac1c\,
    \frac{T''(s)-\{b(s)/a(s)\}\,T'(s)}
         {2T'(s)^2-T(s)T''(s)+\{b(s)/a(s)\}\,T(s)T'(s)}.
\]
Let $T(z)$, $R(z)$ be the meromorphic extensions of $T(s)$, $\overline R(s)$
on an open $D\subset\CC$ containing $I$. Then
\[
  \phi(s,t) = \frac{4\,T_z\,\overline R_{\bar z}}{(1+cT\overline R)^2}
\]
is the unique solution of \eqref{eq:Lcauchy}.
\end{theorem}

For the degenerate case $c=0$ the Cauchy data is instead
$\phi(s,0)=a(s)$, $\phi_t(s,0)=d(s)$ (real-analytic $a>0$, $d$):

\begin{theorem}[{\cite[Thm.\,6]{GMM1}}, case $c=0$]\label{thm:GMM1-6}
Fix $s_0\in I$, put $\theta(s)=-\tfrac12\int_{s_0}^s d(r)/a(r)\,dr$, and
let $\sqrt a(z)$, $\theta(z)$ be the holomorphic extensions of
$\sqrt{a(s)}$, $\theta(s)$. Then $\phi(s,t)=|f(z)|^2$, $z=s+it$, with
\[
  f(z)=\sqrt a(z)\,e^{i\theta(z)}.
\]
\end{theorem}

There is also a purely \emph{geometric} construction of the solution to
the real form of \eqref{eq:Lcauchy},
\begin{equation}\label{eq:Lcauchyreal}
  \Delta(\ln\phi)=-2c\phi, \qquad \phi(s,0)=a(s),\ \phi_t(s,0)=d(s).
\end{equation}
Let $\mathfrak Q(c)$ be the constant-curvature-$c$ Riemann surface
$\mathfrak Q(0)=\RR^2$, $\mathfrak Q(c)=\mathbb S^2(c)$ ($c>0$) or
$\mathbb H^2(c)$ ($c<0$), and $\pi(x_0,x_1,x_2)=\tfrac{x_1+ix_2}{1-cx_0}$
its stereographic projection to $\Cbar$.

\begin{theorem}[{\cite[Thm.\,3]{GMM1}}]\label{thm:GMM1-3}
The unique solution of \eqref{eq:Lcauchyreal} is
\[
  \phi(s,t)=\frac{4|g'(z)|^2}{(1+c|g(z)|^2)^2}, \qquad z=s+it,
\]
where $g(z)$ is the meromorphic extension of $g(s)=\pi(\alpha(s))$, and
$\alpha(s)$ is the unique arclength-parametrized curve in $\mathfrak Q(c)$
with geodesic curvature $\kappa(s)=-d(s)/(2a(s)^{3/2})$ and speed function
$u(s)=\int^s\sqrt{a}(r)\,dr$.
\end{theorem}

Finally, the \emph{modified Liouville equation} used directly in the
Björling construction is
\begin{equation}\label{eq:Lmod}
  4(\ln\rho)_{z\bar z} = -c\,\rho^2\,|f(z)|^2 ,
\end{equation}
$f$ meromorphic: $\rho$ solves \eqref{eq:Lmod} iff $\phi=\rho^2|f|^2$
solves \eqref{eq:Ldef}. For its Cauchy problem
$\rho(s,0)=v(s)$, $\rho_z(s,0)=w(s)$ ($f$ with no real pole):

\begin{corollary}\label{cor:Lmod}
The solution of the Cauchy problem for \eqref{eq:Lmod} is
$\rho=\sqrt\phi/|f|$, with $\phi$ the solution of \eqref{eq:Lcauchy}
(via Theorem~\ref{thm:GMM1-1} or \ref{thm:GMM1-6}) with
\[
  a(s)=v(s)^2|f(s)|^2, \qquad
  b(s)=2v(s)w(s)|f(s)|^2+v(s)^2 f'(s)\overline{f(s)} .
\]
\end{corollary}

% =====================================================================
\section{The Björling problem for BLW surfaces}\label{sec:BJ}
% =====================================================================

\begin{definition}\label{def:analyticcurve}
A curve $\beta:I\to S$ on a Riemann surface $S$ is \emph{analytic at
$s_0\in I$} if, in a conformal chart $(U,z)$ around $\beta(s_0)$, the
curve $z\circ\beta$ is real-analytic at $s_0$; $\beta$ is \emph{analytic}
if this holds at every $s\in I$. Regularity of $\beta$ is defined
analogously.
\end{definition}

\subsection{Statement of the problem}\label{ssec:statement}

\begin{problem}[Björling problem for BLW surfaces]\label{prob:BJ}
Let $a,b\in\RR$ with $|a+b|=1$, let $\beta:I\to\Hyp$ be a regular
analytic curve and $V:I\to\deS$ an analytic vector field along $\beta$
such that, for all $s\in I$,
\begin{align}
  &\ip{\beta(s)}{V(s)} = \ip{\beta'(s)}{V(s)} = 0, \label{eq:cond1}\\
  &\ip{\beta'(s)}{a\beta'(s)-bV'(s)} > 0. \label{eq:cond2}
\end{align}
Find a BLW surface $\psi:S\to\Hyp$ satisfying \eqref{eq:BLWdef} such that
there is a regular analytic curve $\Gamma\subset S$ with $\psi(\Gamma)=\beta$
and with unit normal field of $\psi$ along $\beta$ equal to $V$. Such a
$\psi$ is a \emph{solution of the Björling problem} with \emph{initial
data} $(\beta,V)$.
\end{problem}

Condition \eqref{eq:cond2} is exactly the requirement that the induced
conformal factor $\chi$ of $\sigma=aI+bI\!I$ be positive along $\beta$
(see \eqref{eq:sigma-along-beta} below), i.e.\ that $\sigma$ have a chance
of being positive-definite near $\beta$.

\subsection{Main theorem}

\begin{theorem}[Existence and uniqueness]\label{thm:BJ}
Let $(\beta,V)$ be initial data as in Problem~\ref{prob:BJ}. Suppose
$v=\beta+V$ satisfies $v_0(s)+v_3(s)\neq0$ for all $s\in I$, and that the
map
\begin{equation}\label{eq:Gextra}
  G = \frac{v_1-iv_2}{v_0+v_3} : I\to\CC
\end{equation}
has nowhere-vanishing derivative on $I$. Then the Björling problem with
initial data $(\beta,V)$ has a \emph{unique} solution $\psi:D\subset\CC\to
\Hyp$, constructed near $\beta$ as follows.
\begin{description}
\item[Case $a+b=1$.] Let $\beta(z)$, $V(z)$ be the holomorphic
  extensions of $\beta(s)$, $V(s)$; set $v(z)=\beta(z)+V(z)$ and
  \begin{equation}\label{eq:Gz}
    G(z) = \frac{v_1(z)-iv_2(z)}{v_0(z)+v_3(z)}.
  \end{equation}
  Let $\rho:D\to\CC\setminus\{0\}$ solve the Cauchy problem
  \begin{equation}\label{eq:Lmodcompleted}
    \begin{cases}
      4(\log\rho)_{z\bar z} = -\dfrac{a}{a+b}\,\rho^2|G_z|^2,\\[1mm]
      \rho(s,0) = v_0(s)+v_3(s),\\[1mm]
      \rho_z(s,0) = \dfrac12\Big\{v_0'+v_3'
        + i\big[(V\wedge v')_0+(V\wedge v')_3\big]\Big\};
    \end{cases}
  \end{equation}
  then $\psi = F\Omega F^* : D\to\Hyp$, where
  \begin{equation}\label{eq:FOmega}
    F=\begin{pmatrix}1&0\\ G&1\end{pmatrix}, \qquad
    \Omega = \begin{pmatrix}
      \dfrac{2a+b}{2(a+b)}\rho + \dfrac{2\rho_{z\bar z}}{\rho^2|G_z|^2}
        & \dfrac{2\rho_z}{\rho^2 G_z}\\[3mm]
      \dfrac{2\rho_{\bar z}}{\rho^2\overline{G_z}} & \dfrac2\rho
    \end{pmatrix}.
  \end{equation}
\item[Case $a+b=-1$.] Let $\beta(z)$, $V(z)$ be the \emph{anti}-holomorphic
  extensions of $\beta(s)$, $V(s)$; set $v(z)=\beta(z)+V(z)$ and $G(z)$ as
  in \eqref{eq:Gz}. Let $\rho:D\to\CC\setminus\{0\}$ solve
  \begin{equation}\label{eq:Lmodcompleted2}
    \begin{cases}
      4(\log\rho)_{z\bar z} = -\dfrac a{a+b}\,\rho^2|\overline{G_z}|^2,\\[1mm]
      \rho(s,0) = v_0(s)+v_3(s),\\[1mm]
      \rho_z(s,0) = \dfrac12\Big\{v_0'+v_3'
        - i\big[(V\wedge v')_0+(V\wedge v')_3\big]\Big\};
    \end{cases}
  \end{equation}
  then $\psi=F\Omega F^*:D\to\Hyp$ with $F$ as before and
  \begin{equation}\label{eq:FOmega2}
    \Omega = \begin{pmatrix}
      \dfrac{2a+b}{2(a+b)}\rho + \dfrac{2\rho_{z\bar z}}{\rho^2|\overline{G_z}|^2}
        & \dfrac{2\rho_{\bar z}}{\rho^2 G_{\bar z}}\\[3mm]
      \dfrac{2\rho_z}{\rho^2\overline{G_z}} & \dfrac2\rho
    \end{pmatrix}.
  \end{equation}
\end{description}
In either case, $G$ is the hyperbolic Gauss map of $\psi$, and $\rho$ is
constructed via Section~\ref{ssec:Liouville}.
\end{theorem}

Two remarks on the hypotheses (elaborated in Section~\ref{ssec:remarks}):
the non-vanishing of $v_0+v_3$ can be removed (composing $G$ with a
Möbius transformation), while the non-vanishing of $G'$ genuinely limits
singularities and cannot be dropped for free (Section~\ref{sec:conclusion},
Problems~\ref{open:G.equiv.0}, \ref{open:singularity.1}).

The theorem \emph{contains} the Björling theorems for $\mathtt{CMC1}$
surfaces \cite{GMM1} ($b=0$) and regular flat surfaces \cite{GMM5}
($a=0$) as special cases, and gives a \emph{Björling representation} of
BLW surfaces that bypasses the conformal representation of \cite{GMM2}.

\subsection{Meromorphicity of the hyperbolic Gauss map}

By Bryant \cite{BRYANT}, the hyperbolic Gauss map of a regular surface
$\psi:S\to\Hyp$ is meromorphic if and only if $\psi$ is $\mathtt{CMC1}$.
The following theorem partially generalizes this fact and is the key
technical tool used to justify meromorphic vs.\ anti-meromorphic
extension of $G$ in the two cases of Theorem~\ref{thm:BJ}.

\begin{theorem}\label{thm:G}
Let $\psi:S\to\Hyp$ be a BLW surface satisfying \eqref{eq:BLWdef}, with
$\psi(S)$ not (part of) a horosphere. Then the hyperbolic Gauss map $G$
of $\psi$ is meromorphic if $a+b>0$, and anti-meromorphic if $a+b<0$.
\end{theorem}

For $\mathtt{CMC1}$ surfaces $a+b=a=\pm1$: Bryant's theorem is the special
case $a=1$ (meromorphic); the sign convention in the BLW normalization
absorbs the case $a=-1$ as the anti-meromorphic statement, consistent
with the classical result up to orientation.

\subsection{Sketch of proof}\label{ssec:proofsketch}

The proof of Theorem~\ref{thm:BJ} is organized into a uniqueness part
(Section~\ref{ssec:uniq}) and an existence part
(Section~\ref{ssec:exist}) for the case $a+b=1$; the case $a+b=-1$
(Section~\ref{ssec:case-1}) is handled by an entirely analogous argument.
Section~\ref{ssec:remarks} proves Theorem~\ref{thm:G} and discusses the
hypotheses on the initial data.

\subsubsection{Uniqueness ($a+b=1$)}\label{ssec:uniq}

Suppose $\psi:S\to\Hyp$ solves Problem~\ref{prob:BJ}. Equip $S$ with the
Riemann surface structure induced by $\sigma=aI+bI\!I$ (Lemma~\ref{lem:sigma}).
The strategy is to show that $\psi$ is \emph{locally determined} by the
initial data $(\beta,V)$ via an explicit formula, and then invoke analytic
continuation to get global uniqueness.

\emph{Parametrization.} Choose a conformal chart $(U,z)$ meeting $\Gamma$
in a segment $J'$; write $\gamma=z\circ\Gamma$. Since $\gamma$ is analytic
and regular at a fixed $s_0\in J'$, the inverse function theorem gives a
biholomorphism $\gamma:D\to B$ onto a neighborhood in $z(U)$, and the
surface is represented locally as $\psi(w)=\chi(\gamma(w))$ with
$\psi(s,0)=\beta(s)$, $\eta(s,0)=V(s)$ on $J=(s_0-\delta,s_0+\delta)$.

\emph{Reduction to the modified Liouville equation.} Writing
$v=\beta+V:I\to\LCone$, the hyperbolic Gauss map restricted to $\beta$ is
$G(s)=(v_1-iv_2)/(v_0+v_3)$; since $G$ is meromorphic (justified in
Section~\ref{ssec:remarks}) it extends uniquely to $G(z)$ on $D$. Writing
\begin{equation}\label{eq:psipluseta-rep}
  \psi+\eta = \rho\begin{pmatrix}1&\overline G\\ G&G\overline G\end{pmatrix},
  \qquad \rho=N_0+N_3,\ N=\psi+\eta,
\end{equation}
the fact that $\tau=\ip{d(\psi+\eta)}{d(\psi+\eta)}$ has constant
curvature $c=a/(a+b)$ (Remark~\ref{rmk:psipluseta}) together with
\eqref{eq:psipluseta-rep} gives $\ip{(\psi+\eta)_z}{(\psi+\eta)_{\bar z}}
=\tfrac12\rho^2|G_z|^2$, so the conformal factor $\phi=\rho^2|G_z|^2$ of
$\tau$ satisfies the Liouville equation \eqref{eq:Ldef}; since $G_z$ is
meromorphic, $\rho$ itself satisfies the \emph{modified} Liouville
equation \eqref{eq:Lmod} with $c=a/(a+b)$.

\emph{Initial conditions.} Directly, $\rho(s,0)=v_0(s)+v_3(s)$. To find
$\rho_z(s,0)$, expand $\eta_z,\eta_{\bar z}$ via \eqref{eq:eta-expansion}
to get $(\psi+\eta)_z\wedge(\psi+\eta)_{\bar z}=-(2H-K-1)\psi_z\wedge\psi_{\bar z}$,
and combine with \eqref{eq:2H-K-1} (valid, by Bryant/\cite{GMM1}, even when
$b=0$) to get $\sgn(2H-K-1)=\sgn(a+b)$ at regular points, hence
\[
  (\psi+\eta)_t = -\sgn(a+b)\,\eta\wedge(\psi+\eta)_s.
\]
This yields $(\psi+\eta)_z(s,0)=\tfrac12\{v'+i\sgn(a+b)\,V\wedge v'\}$, i.e.
\begin{equation}\label{eq:rhoz-init}
  \rho_z(s,0) = \frac12\Big\{v_0'+v_3'
    + i\sgn(a+b)\big[(V\wedge v')_0+(V\wedge v')_3\big]\Big\},
\end{equation}
which together with the previous two facts is exactly the Cauchy problem
\eqref{eq:Lmodcompleted} (for $a+b=1$; sign flips to
\eqref{eq:Lmodcompleted2} when $a+b=-1$). Thus $\rho$ is uniquely
determined by $(\beta,V)$.

\emph{Recovering $\psi$.} With $F=\begin{psmallmatrix}1&0\\G&1\end{psmallmatrix}$
one has $\psi+\eta=F\begin{psmallmatrix}\rho&0\\0&0\end{psmallmatrix}F^*$,
and differentiating using $F_z=F\begin{psmallmatrix}0&0\\G_z&0\end{psmallmatrix}$
gives explicit expressions for $(\psi+\eta)_{\bar z}$ and
$(\psi+\eta)_{z\bar z}$ in terms of $F$ and $\rho$. Substituting into the
recovery formula \eqref{eq:psi-recovered} from Remark~\ref{rmk:psipluseta}
gives, at regular points of $\psi+\eta$,
\[
  \psi = F\Omega F^*
\]
with $\Omega$ exactly \eqref{eq:FOmega}, and $\det\Omega=1$ because $\rho$
solves \eqref{eq:Lmodcompleted}. This is the desired closed form,
completing the uniqueness argument (via analytic continuation of BLW
surfaces).

\subsubsection{Existence ($a+b=1$)}\label{ssec:exist}

Given data $(\beta,V)$ satisfying the hypotheses of Theorem~\ref{thm:BJ}
with $a+b=1$: define $G(s)$ by \eqref{eq:Gextra}, extend it meromorphically
to $G(z)$ by \eqref{eq:Gz} on a domain $D$
containing $I$, and let $\rho$ solve \eqref{eq:Lmodcompleted}. On
$U=D\setminus P$ (removing the closed set $P$ of zeros of $\rho$ and
zeros/poles of $G_z$ — a set that does not meet $I$), define
$\psi=F\Omega F^*$ and $N=F\begin{psmallmatrix}\rho&0\\0&0\end{psmallmatrix}F^*$
with $F,\Omega$ as in \eqref{eq:FOmega}. The proof that $\psi$ is a
solution proceeds in five steps.

\begin{enumerate}[label=(\roman*)]
\item \textbf{$\sigma=aI+bI\!I$ is positive-definite.} Writing
  $\Omega=\begin{psmallmatrix}A&B\\ \overline B&C\end{psmallmatrix}$ and
  $U:=AG_z+\overline B_z=\tfrac{b}{2(a+b)}\rho G_z$ (using
  \eqref{eq:Lmodcompleted}), direct computation with
  $F_z=F\begin{psmallmatrix}0&0\\G_z&0\end{psmallmatrix}$ shows the
  $dz^2$- and $d\bar z^2$-coefficients of $\sigma$ vanish and the
  remaining ($|dz|^2$) coefficient is
  \begin{equation}\label{eq:sigma-expanded}
    2\chi = (a+b)\big\{|B_z|^2-|U|^2\big\}.
  \end{equation}
  Along $\beta$ this equals
  \begin{equation}\label{eq:sigma-along-beta}
    2\chi(s,0) = 2\ip{\beta'(s)}{a\beta'(s)-bV'(s)} > 0
  \end{equation}
  by hypothesis \eqref{eq:cond2}, so $\sigma$ is positive-definite (hence
  a genuine conformal metric) near $\beta$.
\item \textbf{$\psi$ is regular.} From
  $\psi_z=F\begin{psmallmatrix}*&B_z\\U&0\end{psmallmatrix}F^*$ one
  computes $\ip{\psi_s\wedge\psi_t}{\psi_s\wedge\psi_t}=|B_z|^2-|U|^2$,
  which is exactly the (positive, by step (i)) coefficient of $\sigma$;
  hence $\psi_s\wedge\psi_t\neq0$ near $\beta$.
\item \textbf{$\eta:=N-\psi$ is the unit normal.} Using $\det\Omega=1$
  and $\rho C=2$ one checks $\ip\eta\eta=1$ and $\ip\psi\eta=0$ by direct
  matrix-determinant identities; finally
  $\ip{\psi_z\wedge\psi_{\bar z}}\eta = \tfrac i2\{|B_z|^2-|U|^2\}$ (computed
  via the symbolic routines of Appendix~B), and since $a+b=1$,
  \eqref{eq:sigma-expanded} forces $\ip{\psi_s\wedge\psi_t}\eta>0$, i.e.\
  $\eta$ has the correct orientation.
\item \textbf{$\psi$ is BLW.} Substituting the computed $d\eta$ into
  \eqref{eq:eta-expansion} gives a $2\times2$ linear system for $H,\Upsilon$
  with nonzero determinant $D=|U|^2-|B_z|^2$, whose solution is
  \[
    H = 1-\frac{b}{2(a+b)}\frac{\rho^2|G_z|^2}{D}, \qquad
    K = H^2-|\Upsilon|^2 = 1+\frac{a}{a+b}\frac{\rho^2|G_z|^2}{D},
  \]
  from which $2a(H-1)+b(K-1)=0$ holds identically.
\item \textbf{$\psi$ realizes the data $(\beta,V)$.} An identity relating
  the initial curve data,
  \begin{equation}\label{eq:init-prop}
    (v_0+v_3)^2 G'\big\{\beta_1+i\beta_2-\overline G(\beta_0+\beta_3)\big\}
      = v_0'+v_3'+i\big[(V\wedge v')_0+(V\wedge v')_3\big],
  \end{equation}
  combined with the third line of \eqref{eq:Lmodcompleted}, gives
  $B(s,0)=\beta_1+i\beta_2-\overline G(\beta_0+\beta_3)$, and then
  ($\det\Omega=1$, $\rho C=2$) $A(s,0)=\beta_0+\beta_3$. Writing out
  $\psi=F\Omega F^*$ entrywise along $\beta$ and using
  $\ip\beta\beta=-1$ (so $\beta_0+\beta_3\neq0$) pins down the remaining
  matrix entry, giving $\psi(s,0)=\beta(s)$ for all $s\in I$; since
  $(\psi+\eta)(s,0)=v(s)=\beta(s)+V(s)$ by construction, this also gives
  $\eta(s,0)=V(s)$.
\end{enumerate}
Together, (i)–(v) show $\psi$ solves the Björling problem with data
$(\beta,V)$, completing the existence proof for $a+b=1$.

\subsubsection{The case $a+b=-1$}\label{ssec:case-1}

The argument mirrors Sections~\ref{ssec:uniq}--\ref{ssec:exist} verbatim,
with $G$ now extended \emph{anti}-meromorphically (justified by
Theorem~\ref{thm:G} for $a+b<0$). The representation
\eqref{eq:psipluseta-rep} for $\psi+\eta$ is unchanged, but now
$\ip{(\psi+\eta)_z}{(\psi+\eta)_{\bar z}}=\tfrac12\rho^2|\overline{G_z}|^2$,
giving the modified Liouville equation with $\overline{G_z}$ in place of
$G_z$, and Cauchy data \eqref{eq:Lmodcompleted2} (a sign flip in the
imaginary part relative to \eqref{eq:Lmodcompleted}, coming from
$\sgn(a+b)=-1$ in \eqref{eq:rhoz-init}). The recovery formula gives
$\psi=F\Omega F^*$ with $\Omega$ as in \eqref{eq:FOmega2}. All five
existence steps go through with the roles of $U,B_z$ swapped
appropriately, giving
\[
  H = 1-\frac{b}{2(a+b)}\frac{\rho^2|\overline{G_z}|^2}{D}, \qquad
  K = 1+\frac{a}{a+b}\frac{\rho^2|\overline{G_z}|^2}{D}, \qquad
  D=|U|^2-|\overline{B_z}|^2\neq0,
\]
and $\psi(s,0)=\beta(s)$ follows exactly as before.

\subsection{Further remarks}\label{ssec:remarks}

\subsubsection{On the hypotheses on the initial data}

Writing $Q=a\psi-b\eta$ and using $\sigma=\ip{d\psi}{dQ}=2\chi|dz|^2$
in a $\sigma$-conformal chart, one gets $2\chi(s,0)=
\ip{\beta'(s)}{a\beta'(s)-bV'(s)}>0$: this is exactly why condition
\eqref{eq:cond2} is imposed in the statement of the problem — it is a
\emph{necessary} condition for any solution to exist, not an artificial
restriction.

The hypothesis $v_0(s)+v_3(s)\neq0$ (needed for \eqref{eq:Gz} to make
sense) can be dropped following \cite{GMM1}: composing $G$ with a
suitable Möbius transformation of $\Cbar$ reduces to this case. Concretely,
since $v(s)=\beta(s)+V(s)\neq0$, whenever $v_0+v_3=0$ one has
$v_0-v_3\neq0$, and one instead writes
$N=\psi+\eta=\rho_*\begin{psmallmatrix}H\overline H&\overline H\\H&1\end{psmallmatrix}$,
$\rho_*=N_0-N_3$, $H=(N_1-iN_2)/(N_0-N_3)$, with the extra requirement
$H'(s_0)\neq0$ at the exceptional point; the same computations then give
the same conclusion of Theorem~\ref{thm:BJ} by locality of the problem.

The nowhere-vanishing of $G'$ is more essential. If $G'\equiv0$
identically, Remark~\ref{rmk:psipluseta} forces any solution to be (part
of) a horosphere. If $G'\not\equiv0$ but has isolated zeros (forced to be
isolated by analyticity), the construction of Theorem~\ref{thm:BJ} can
develop singularities; analyzing these singularities is left as an open
problem (Section~\ref{sec:conclusion}, Problems~\ref{open:G.equiv.0},
\ref{open:singularity.1}).

\subsubsection{Why $G$ is meromorphic: proof of Theorem~\ref{thm:G}}

By Remark~\ref{rmk:psipluseta}(1), $G$ is always conformal, hence
meromorphic or anti-meromorphic; the content of Theorem~\ref{thm:G} is
pinning down \emph{which one}, by the sign of $a+b$. The uniqueness
argument of Section~\ref{ssec:uniq} tacitly assumes $G$ meromorphic for
$a+b=1$; this is justified by contradiction: assume $\psi$ is BLW,
$a+b=1$, $\sigma$ positive-definite, $\psi(S)$ not a horosphere (so the
umbilic points — equivalently the zeros of $G_{\bar z}$ — are isolated),
and suppose \emph{towards contradiction} that $G$ is anti-meromorphic.
Fix a conformal chart $(U,z)$ (preserving orientation) and write
$\psi+\eta$ as in \eqref{eq:psipluseta-rep} with $\rho$ nonvanishing on
$U$ (achievable after composing $G$ with a suitable Möbius map). The
computation of Section~\ref{ssec:case-1} (the $a+b=-1$ formulas, which
apply here because they only used anti-meromorphicity of $G$, not the
sign of $a+b$) gives $\psi=F\Omega F^*$ with $\Omega$ as in
\eqref{eq:FOmega2}, hence \eqref{eq:sigma-expanded} for $a+b=1$ forces
$|\overline{B_z}|^2-|U|^2>0$; but then the analogue of the identity used in
step (iii) of Section~\ref{ssec:exist} gives
$\ip{\psi_s\wedge\psi_t}\eta<0$ — contradicting that $\eta$ and
$\psi_s\wedge\psi_t$ point in the same direction (by construction of the
unit normal, Section~\ref{ssec:local}). Hence $G$ must be meromorphic
when $a+b=1$; the case $a+b=-1$ (anti-meromorphic) is symmetric.

% =====================================================================
\section{Conclusions and open problems}\label{sec:conclusion}
% =====================================================================

The thesis poses and solves the Björling problem for BLW surfaces.
Theorem~\ref{thm:BJ} establishes existence and uniqueness of the solution;
its representation formula does not use the conformal representation of
\cite{GMM2}, and is thus a genuine \emph{Björling representation} for BLW
surfaces. It specializes to the known results for $\mathtt{CMC1}$ surfaces
\cite{GMM1} and regular flat surfaces \cite{GMM5}. Theorem~\ref{thm:G}
determines the meromorphicity type of the hyperbolic Gauss map of a BLW
surface from the sign of $a+b$. Along the way, the author establishes the
Lorentzian cross-product identities \eqref{eq:wpbase}--\eqref{eq:wpb},
which are used repeatedly and are noted as small original computational
results.

For lack of time, the thesis does not address one of the important
consequences of solving the Björling problem: the \emph{symmetry
principle} for BLW surfaces. This should follow by an argument entirely
parallel to \cite{GMM1} once an appropriate notion of symmetry (for BLW
surfaces reflected across geodesic planes orthogonal to them) is set up.

\subsection*{Open problems}
The following related problems (some already studied for $\mathtt{CMC1}$
surfaces in \cite{GMM1}, as noted) are proposed:
\begin{enumerate}[label=\arabic*.]
\item Study BLW surfaces with the topology of a cylinder.
\item Study complete BLW surfaces with special finite \emph{dual} total
  curvature.
\item Find an explicit closed formula for the solution when the initial
  curve $\beta$ is planar.
\item\label{open:G.equiv.0} Solve the Björling problem when $G$ in
  \eqref{eq:Gextra} is a \emph{constant} map; as noted in
  Section~\ref{ssec:remarks}, any solution in this case must be (part of)
  a horosphere.
\item Find a geometric solution of the Cauchy problem for the Liouville
  equation \eqref{eq:Lcauchyreal} (with $c=1$) directly, in the spirit of
  \cite{GMM1}'s geometric approach via the secondary Gauss map (related
  to the rigidity property) of $\mathtt{CMC1}$ surfaces, recovering the
  solution of \eqref{eq:Lcauchyreal} from that map and the initial data.
\item\label{open:singularity.1} Solve the Björling problem allowing
  singularities — e.g.\ where $G'$ or $\ip{\beta'}{a\beta'-bV'}$ vanish at
  isolated points of $I$ — extending the notion of BLW surface to points
  where \eqref{eq:BLWdef} only holds at regular points. This parallels
  \cite{GMM5}'s treatment of flat fronts (which allow singularities,
  unlike flat surfaces) and its subsequent classification of flat fronts
  with embedded isolated singularities.
\item Solve the Björling problem when the initial data satisfies
  $\ip{\beta'}{a\beta'-bV'}\equiv0$ identically (note: by analyticity of
  $\beta,V$, the zero set of this quantity is either isolated or all of
  $I$). By Corollary~10 of \cite{GMM5} (case $a=0$), given $\beta$ and a
  suitably chosen $V$, there is a unique flat front $\Sigma$ for which
  $\beta$ consists exactly of the singular points of $\Sigma$.
\item Solve the Björling problem for $\mathtt{ELW}$ (Elliptic Linear
  Weingarten) surfaces \cite{GMM6} — surfaces in $\RR^3$ satisfying an
  analogous relation between $H$ and $K$. This case may in fact be
  \emph{harder} than the BLW case: for $\mathtt{ELW}$ surfaces only a
  \emph{harmonic}, not conformal, representation is available
  \cite{GMM6}, which is the chief obstacle.
\item\label{open:G} By \cite{BRYANT,GMM1}, the hyperbolic Gauss map of a
  regular surface $\psi:S\to\Hyp$ is meromorphic iff $\psi$ is
  $\mathtt{CMC1}$; Theorem~\ref{thm:G} partially generalizes this. The
  converse direction is open: if $\sigma=aI+bI\!I$ is positive-definite
  and $G$ is meromorphic (or anti-meromorphic), must $\psi$ be BLW?
\end{enumerate}
Finally, there is also a Björling problem for \emph{maximal} surfaces in
$\mathbb L^n$; preliminary computations (Assoc.\ Prof.\ Đoàn Thế Hiếu,
based on \cite{ACM,AM}) suggest this can be solved by integrating
functions with no real periods.

% =====================================================================
\appendix
\section{Previously solved Björling-type problems}\label{app:list}
% =====================================================================

The following table (reproduced from the thesis's Appendix~A) lists
Björling-type problems solved prior to this thesis. In every case the
solution is unique and yields a symmetry principle for the surface class
considered.

\begin{center}
\begin{tabular}{lll c}
\textit{Surface class} & \textit{Reference} & \textit{Author(s)} & \textit{Year}\\
\hline
minimal, $\RR^3$ & \cite{FUNG} & E.\,G.\ Björling, H.\,A.\ Schwarz & 1844, 1890\\
minimal, $\RR^n$ & \cite{AM} & L.\,J.\ Alías, P.\ Mira & 2002\\
maximal, $\mathbb L^3$ & \cite{ACM} & L.\,J.\ Alías et al. & 2003\\
minimal Möbius strip, $\RR^n$ & \cite{MIRA} & P.\ Mira & 2005\\
flat front, $\Hyp$ & \cite{GMM5} & J.\,A.\ Gálvez, P.\ Mira & 2005\\
flat surface, $\Hyp$ & \cite{GMM5} & J.\,A.\ Gálvez, P.\ Mira & 2005\\
$\mathtt{CMC1}$, $\Hyp$ & \cite{GMM1} & J.\,A.\ Gálvez, P.\ Mira & 2005\\
improper affine sphere, $\RR^3$ & \cite{ACG} & J.\,A.\ Aledo et al. & (preprint)\\
\textbf{BLW}, $\Hyp$ & \textit{this thesis} & Huỳnh Kỳ Anh & 2006\\
\end{tabular}
\end{center}

% =====================================================================
\section{Note on the appendix on symbolic computation}
% =====================================================================

The thesis includes an appendix with Maple routines used to verify, by
symbolic (Hermitian $2\times2$ matrix) computation, identities such as
$\ip{\psi_z\wedge\psi_{\bar z}}\eta = \tfrac i2\{|B_z|^2-|U|^2\}$ (used in
step (iii) of Section~\ref{ssec:exist}) for the
case $a+b=1$, and the analogous identity with $U,\overline{B_z}$ swapped
for $a+b=-1$ — that are otherwise long, purely mechanical determinant
expansions. These routines are a computational aid, not a separate
mathematical contribution, and are not reproduced here.

% =====================================================================
\begin{thebibliography}{99}\small

\bibitem{GMM1} J.\,A.\ Gálvez, A.\ Martínez, P.\ Mira, \emph{The Björling
  problem for constant mean curvature one surfaces in the hyperbolic
  space}, related works of Gálvez--Mira on $\mathtt{CMC1}$ surfaces.

\bibitem{GMM2} J.\,A.\ Gálvez, A.\ Martínez, P.\ Mira, \emph{On the
  conformal (harmonic) representation of linear Weingarten surfaces of
  Bryant type in $\Hyp$} (referenced as a preprint at the time of
  writing).

\bibitem{GMM5} J.\,A.\ Gálvez, P.\ Mira, \emph{Embedded isolated
  singularities of flat surfaces in hyperbolic space} and the Björling
  problem for flat fronts/flat surfaces, 2005.

\bibitem{GMM6} J.\,A.\ Gálvez, A.\ Martínez, P.\ Mira, on Elliptic Linear
  Weingarten ($\mathtt{ELW}$) surfaces in $\RR^3$.

\bibitem{BRYANT} R.\ Bryant, \emph{Surfaces of mean curvature one in
  hyperbolic space}, Astérisque \textbf{154--155} (1987).

\bibitem{MAGIC} Reference on the degenerate BLW case $a+b=0$ (cited via
  \cite{GMM2}; not available to the author of the thesis).

\bibitem{FKR} Local theory of surfaces reference (Fujimori--Kobayashi--Rossman
  or similar).

\bibitem{FUJIMORI} S.\ Fujimori, concise treatment of the basic equations
  of surface theory.

\bibitem{CARMO} M.\ P.\ do Carmo, treatment of the hyperbolic space model
  with detailed proofs.

\bibitem{FUNG} K.\ Fung, treatment of minimal surfaces in $\RR^3$ as
  isotropic curves in $\CC^3$, with a fast solution of the classical
  Björling problem.

\bibitem{AM} L.\,J.\ Alías, P.\ Mira, Björling problem for minimal
  surfaces in $\RR^n$, 2002.

\bibitem{ACM} L.\,J.\ Alías et al., Björling problem for maximal surfaces
  in $\mathbb L^3$, 2003.

\bibitem{MIRA} P.\ Mira, Björling problem for minimal Möbius strips in
  $\RR^n$, 2005.

\bibitem{ACG} J.\,A.\ Aledo et al., Björling problem for improper affine
  spheres in $\RR^3$ (preprint).

\bibitem{LIOUVILLE} Reference(s) on the Liouville equation and its Cauchy
  problem.

\bibitem{HITCHIN} N.\ Hitchin, lecture notes on Riemann surfaces.

\bibitem{WIKI:RS} Wikipedia, \emph{Riemann surface}.

\bibitem{WIKI:AC} Wikipedia, \emph{Analytic continuation}.

\bibitem{AHLFORS} L.\ Ahlfors, \emph{Complex Analysis}.

\end{thebibliography}

\end{document}
